\begin{resumo}
\begin{SingleSpace}
A pandemia de Covid-19 evidenciou a necessidade de estratégias logísticas ágeis para evitar o desperdício de imunizantes e salvar vidas. Diante disso, este trabalho objetiva minimizar o tempo de distribuição e mitigar a perda de vacinas por expiração de validade, utilizando modelos de Programação Linear Inteira (PLI) aplicados aos níveis intermunicipal e intramunicipal. A metodologia baseia-se no desenvolvimento de dois modelos matemáticos integrados: um de redistribuição direta entre municípios e outro que considera a intermediação de centros estaduais. As formulações incorporam restrições operacionais complexas, como múltiplas origens e destinos, diferentes modais de transporte (aéreo e rodoviário), capacidades de carga agregada e penalidades baseadas na urgência biológica dos lotes. A resolução computacional foi implementada na linguagem Python com o auxílio do solver Gurobi Optimizer, avaliando um estudo de caso nacional composto pelas 27 capitais brasileiras, alguns dos seus respectivos postos de saúde e quatro tipos de imunizantes. Os resultados demonstram a eficácia da abordagem exata frente à alta complexidade combinatória do problema. A análise dos cenários, complementada por visualizações geográficas interativas, revela que a introdução da penalidade de perecibilidade força o sistema a reconfigurar sua malha estrutural. O solver prioriza o uso de rotas aéreas de longa distância para escoar estoques críticos antes de sua expiração, sobrepondo a urgência temporal à proximidade geográfica. Conclui-se que a formulação matemática proposta, aliada ao processamento computacional avançado, configura uma ferramenta tática viável e escalável para o enfrentamento de crises sanitárias. A transição de um planejamento empírico para um sistema automatizado fundamentado em Pesquisa Operacional permite a alocação ótima de recursos e consolida um legado tecnológico aplicável à gestão da saúde pública.
\end{SingleSpace}
\vspace{\onelineskip}
\textbf{Palavras-chave}: logística de vacinas; programação linear inteira; pesquisa operacional; otimização; distribuição de vacinas; Gurobi.
\end{resumo}